\documentclass[11pt,a4paper]{article}

%% PACHETE MATEMATICE
\usepackage{amsmath,amssymb,amsfonts}
\usepackage{amsthm}
\usepackage{mathpazo}
\usepackage{eulervm}
\usepackage{enumerate}
\usepackage{color}
\usepackage{graphicx}
\usepackage[all]{xy}
\usepackage{xcolor}
\usepackage{framed}
\usepackage[unicode]{hyperref}
\setcounter{MaxMatrixCols}{20}
\usepackage{multicol}
%\usepackage[romanian]{babel}


%% ASPECT
\linespread{1.05}
\usepackage[T1]{fontenc}
\usepackage[utf8x]{inputenc}
\usepackage[margin=2cm]{geometry}
% \usepackage[color]{showkeys}
\usepackage{anyfontsize}
\usepackage{titlesec}
\titleformat*{\section}{\large\bfseries}

\usepackage{fancyhdr}
\pagestyle{fancy}
\rhead{\small \color{gray} Notițe de Adrian Manea}
\lhead{\small \color{gray} ETTI-AN1, 2017---2018, C. Ghiu}


\usepackage[autostyle = false, style = german]{csquotes}
\usepackage[romanian]{babel}
\newcommand\qq{\enquote}

%% COMENZILE MELE
\renewcommand*{\proofname}{Dem.:}
\newcommand\bb{\mathbb}
\newcommand\dr{\mathrm}
\newcommand\kal{\mathcal}
\newcommand\fk{\mathfrak}
\newcommand\op{\oplus}
\newcommand\ot{\otimes}
\newcommand\ds{\displaystyle}
\newcommand\id{\indent}
\newcommand\nid{\noindent}
\newcommand\seq{\subseteq}
\newcommand\sneq{\subsetneq}
\newcommand\speq{\supseteq}
\newcommand\spneq{\supsetneq}
\newcommand\rar{\rightarrow}
\newcommand\Rar{\Rightarrow}
\newcommand\xrar{\xrightarrow}
\newcommand\lar{\leftarrow}
\newcommand\Lar{\Leftarrow}
\newcommand\xlar{\xleftarrow}
\newcommand\lrar{\leftrightarrow}
\newcommand\Lrar{\Leftrightarrow}
\newcommand\ideal{\trianglelefteq}
\newcommand\ai{\text{ a.\^{i}. }}
\newcommand\sk{(\textit{Schi\c{t}\u{a}}) }
\newcommand\rhup{\rightharpoonup}
\newcommand\rhdn{\rightharpoondown}
\newcommand\lhup{\leftharpoonup}
\newcommand\lhdn{\leftharpoondown}
\newcommand\vid{\emptyset}
\newcommand\wh{\widehat}
\newcommand\ol{\overline}
\newcommand\NN{\mathbb{N}}
\newcommand\RR{\mathbb{R}}
\newcommand\ZZ{\mathbb{Z}}
\newcommand\QQ{\mathbb{Q}}
\newcommand\CC{\mathbb{C}}

%% STILURILE MELE
\newtheoremstyle{dotless}{}{}{\itshape}{}{\bfseries}{:}{ }{}

\theoremstyle{dotless}
\newtheorem{theorem}{Teorem\u{a}}[section]
\newtheorem{proposition}{Propozi\c{t}ie}[section]
\newtheorem{lemma}{Lem\u{a}}[section]
\newtheorem{corollary}{Corolar}[section]
\newtheorem{exercise}{Exerci\c{t}iu}

\newtheoremstyle{dotlessdr}{}{}{}{}{\bfseries}{:}{ }{}
\theoremstyle{dotlessdr}
\newtheorem{example}{Exemplu}[section]
\newtheorem{definition}{Defini\c{t}ie}[section]
\newtheorem{remark}{Observa\c{t}ie}[section]




\begin{document}
% \definecolor{labelkey}{RGB}{222,0,0}

\begin{center}
  {\large\textbf{Seminar 3 \\ Exerciții spații vectoriale}}
\end{center}

\vspace{1cm}

1. Fie $S = \Bigg\{ \begin{pmatrix}
  0 & a + b & 0 \\
  a & 0 & b
\end{pmatrix} \mid a, b \in \RR \Bigg\}$.

Arătați că $S$ este subspațiu vectorial al lui $M_{2,3}(\RR)$.

\vspace{1cm}

2. Decideți care din următoarele sisteme de vectori sînt liniar independente:
\begin{enumerate}[(a)]
\item $S_1 = \Bigg\{ v_1 = \begin{pmatrix} 1 & 2 \\ 1 & 1 \end{pmatrix}, %
  v_2 = \begin{pmatrix} 2 & -1 \\ 1 & 3 \end{pmatrix}, %
  v_3 = \begin{pmatrix} 3 & -2 \\ 1 & 1 \end{pmatrix} \Bigg\}$ în $M_2(\RR)$;
\item $S_2 = \Bigg\{ v_1 = \begin{pmatrix} 1 \\ 2 \\ -1 \\ 3 \end{pmatrix}, %
  v_2 = \begin{pmatrix} 4 \\ -1 \\ 2 \\ 1 \end{pmatrix}, %
  v_3 = \begin{pmatrix} 2 \\ 1 \\ 1 \\ -1 \end{pmatrix}, %
  v_4 = \begin{pmatrix} 3 \\ 0 \\ 0 \\ 5 \end{pmatrix} \Bigg\}$ în $\RR^4$;
\item $S_3 = \{ v_1 = -X^2 + X + 1, v_2 = X^2 - X + 1, v_3 = X^2 + X - 1 \}$ în $\RR[X]_{\leq 2}$.
\end{enumerate}

\vspace{1cm}

3. Fie $S = \Bigg\{ v_1 = \begin{pmatrix} 1 & 2 \\ 2 & \alpha \end{pmatrix}, %
v_2 = \begin{pmatrix} 1 & 0 \\ \beta & 1 \end{pmatrix}, %
v_3 = \begin{pmatrix} -1 & 1 \\ 1 & 0 \end{pmatrix} \Bigg\}$.

Aflați $\alpha, \beta \in \RR$, astfel încît sistemul să fie liniar independent peste $\RR$.

\vspace{1cm}

4. Fie $S = \{ v_1 = X^3 + X + 2, v_2 = 2X^3 - X^2 + 1, v_3 = 2x^2 + 3 \}$.
\begin{enumerate}[(a)]
\item Arătați că $S$ este sistem liniar independent în $\RR[X]_{\leq 3}$;
\item Precizați care dintre polinoamele $P = 4X + 3$ și $Q = X + 1$ aparține spațiului $\mathrm{Sp}(S)$;
\item Este $S$ sistem de generatori pentru $\RR[X]_{\leq 3}$?
\end{enumerate}

\vspace{1cm}

5. Fie $V = \Bigg\{
\begin{pmatrix}
  a & b & c \\
  0 & a + b & a \\
  c & 0 & a
\end{pmatrix} \mid a, b, c \in \RR \Bigg\}$.

\begin{enumerate}[(a)]
\item Arătați că $V$ este subspațiu vectorial al lui $M_3(\RR)$.
\item Determinați dimensiunea lui $V$ și precizați o bază a acestui spațiu.
\end{enumerate}

\vspace{1cm}

6. Fie aplicația:
\[
  f : \RR^3 \to \RR^4, \quad f(x_1, x_2, x_3) = (3x_1 + 2x_2, x_2 + x_3, x_1 + x_2, x_1 + 2x_2 - x_3).
\]
\begin{enumerate}[(a)]
\item Arătați că $f$ este aplicație liniară;
\item Determinați $\mathrm{Ker}f$ și $\mathrm{Im}f$;
\item Determinați $\dim \mathrm{Ker} f$ și $\dim \mathrm{Im} f$, precizînd și cîte o
  bază în aceste subspații.
\end{enumerate}




\end{document}

